\section*{Abstract}
\addcontentsline{toc}{section}{Abstract}

University course timetabling is an NP-hard combinatorial optimisation problem in which classes,
instructors, student cohorts and rooms must be assigned to time slots subject to a large set of
hard and soft constraints. The Semester~VI phase of this project established that UniTime, an
open-source constraint-driven academic scheduling system, is able to express every scheduling
constraint identified at COEP Technological University without relaxation or external
post-processing. That phase was a feasibility study: it configured a local deployment and entered
representative test data, but it did not produce a timetable from institutional data.

The present phase, carried out under the project name \emph{taasika-v3}, converted that feasibility
result into a working system. A reproducible data pipeline was built that reads a snapshot of the
university's existing Taasika scheduling database and emits the sixteen ordered XML documents that
UniTime's Data Exchange interface requires, covering the academic session, curriculum structure,
rooms, instructors, course offerings, scheduling preferences and a simulated student body of
1{,}130 students with 9{,}148 section enrolments. Institution-specific rules---minor-degree and
open-elective time blocks, division and batch structure, and dynamic lunch breaks for both students
and instructors---were encoded as generated preference data rather than as manual edits. An
independent evaluation framework was written to score any exported timetable against the same input
data, because the solver's own quality report cannot be used to audit the solver.

Using this framework, the project established a result that the solver's internal counters had been
concealing. Several earlier runs reported zero student conflicts; each was traced to an artefact
rather than to a better timetable, the most common being UniTime silently re-sectioning students
between laboratory batches. Once student enrolments were pinned---by selecting
\emph{Enrolled Student Course Demands} and disabling student switching and re-sectioning---the final
solution scheduled 290 of 291 classes with no instructor and no room double-booking, a hard
feasibility score of 98.85 out of 100, and 40 of 1{,}130 students (3.54\,\%) carrying a timetable
clash, all of them first-year students in a single laboratory--lecture pairing.

A reinforcement-learning layer was also implemented, in which a policy selects among eight stock
CPSolver neighbourhood operators inside the existing search loop. The layer is complete and
instrumented, comprising 25 Java classes and 11 Python modules, and it never writes a placement
itself, so hard-constraint safety is unchanged. Its pre-registered control experiment, however,
answered the project's own research question in the negative: simply raising the solver's
student-conflict weight reduced the violation rate from 0.0414 to below 0.004 on the same instance,
clearing the 25\,\% improvement threshold that the learned policy was intended to achieve. No claim
is therefore made that the learned policy outperforms the baselines; the report states precisely
what was implemented, what was measured, and what remains a hypothesis.
